% ============================================================================
%  SECCIÓN 6.6: ACTUALIZACIÓN DEL REPOSITORIO, EL README Y DIFICULTADES
% ============================================================================

\section{Actualizaci\'on del repositorio, el README y las dificultades}

Al finalizar la integraci\'on se registr\'o el avance en el repositorio,
se document\'o el backend en el README del proyecto y se registraron las
dificultades encontradas durante el desarrollo.

\subsection{Registro del avance en GitHub}

Los cambios del avance se registraron con el flujo habitual de Git:

\begin{lstlisting}[style=terminal]
$ git add .
$ git commit -m "feat(supabase): integracion con Supabase y operaciones CRUD"
$ git push origin docs
\end{lstlisting}

El historial del repositorio evidencia la evoluci\'on del proyecto a
trav\'es de los avances:

\begin{lstlisting}[style=terminal]
Commit 1: inicializacion del proyecto
Commit 2: interfaces y navegacion
Commit 3: logica y validaciones
Commit 4: persistencia local
Commit 5: Supabase + CRUD
\end{lstlisting}

\subsection{Documentaci\'on del backend en el README}

El README incorpor\'o una secci\'on dedicada al backend y a la base de
datos, describiendo la plataforma utilizada, las tablas creadas, la
funcionalidad conectada y las operaciones CRUD implementadas. Adem\'as
incluye la gu\'ia de puesta en marcha del entorno local:

\begin{lstlisting}
## Configuracion de Supabase (Desarrollo Local)

Supabase corre localmente via Docker. Esto levanta:
- PostgreSQL en puerto 54322
- API (PostgREST) en puerto 54321
- Studio (dashboard web) en puerto 54323

### Levantar Supabase

# Iniciar Supabase local
pnpm supabase start

# Resetear DB (migraciones + seed)
pnpm supabase db reset

## Supabase Cloud (Produccion)

# Link local al proyecto cloud
pnpm supabase link --project-ref TU-PROJECT-REF
pnpm supabase db push --seed
\end{lstlisting}

El resultado final puede resumirse as\'i: la aplicaci\'on m\'ovil
desarrollada con Expo se encuentra integrada con Supabase como backend y
utiliza PostgreSQL para la persistencia de la informaci\'on. Se
implementaron operaciones CRUD sobre las entidades principales del
proyecto.

\subsection{Dificultades encontradas}

Durante la integraci\'on se presentaron las siguientes dificultades, que
fueron analizadas y resueltas:

\begin{table}[H]
  \centering
  \caption{Dificultades encontradas en el avance 6}
  \contenidoConFuente{%
    \begin{tablaAPA}
      \begin{tabular}{@{}p{5.2cm}p{4.6cm}p{4.8cm}@{}}
        \toprule
        \textbf{Dificultad} & \textbf{An\'alisis} & \textbf{Soluci\'on} \\
        \midrule
        La aplicaci\'on no recuperaba los
        registros de Supabase & Se revis\'o la configuraci\'on de
        conexi\'on y la consulta realizada &
        Se corrigi\'o la configuraci\'on y se
        verific\'o la consulta nuevamente. \\
        \addlinespace
        Error de recursi\'on infinita en
        las pol\'iticas RLS &
        Las pol\'iticas consultaban la tabla
        \texttt{profiles} de forma recursiva &
        Se cre\'o una funci\'on
        \texttt{get\_user\_role()} con
        \texttt{SECURITY DEFINER} (migraci\'on 004). \\
        \addlinespace
        Los cambios locales no aparec\'ian
        en el proyecto en la nube &
        Las migraciones no se hab\'ian
        sincronizado &
        Se ejecut\'o \texttt{db push --seed}
        para sincronizar esquema y datos. \\
        \addlinespace
        Im\'agenes que no se sub\'ian &
        El bucket \texttt{scenario-images}
        carec\'ia de permisos &
        Se restringi\'o el acceso al rol
        adecuado y se verific\'o la subida. \\
        \bottomrule
      \end{tabular}
    \end{tablaAPA}%
  }{Elaboraci\'on propia en base al registro del equipo de desarrollo.}
\end{table}

Estas dificultades permitieron consolidar el conocimiento sobre la
integraci\'on de una aplicaci\'on m\'ovil con un backend en la nube.